% ===========================================================================
\section{Discussion}
\label{sec:discussion}
% ===========================================================================

\subsection{Choosing the unit is choosing the uncertainty}

The coal blend is unknown, and how much that matters is entirely a function of the unit the
activity data are read in. Per tonne it is a factor of two; per terajoule it is
\SI{6.8}{\percent}. Nothing about Pakistan's coal fleet changed between those two statements.

The general form is worth stating because it recurs. When an activity datum is available in
more than one unit, the right one is the unit in which the unknown property does the least
work---and calorific value is exactly the property that distinguishes the two coals, so
reading energy rather than mass removes it from the calculation rather than requiring it to be
assumed.

This has a corollary that cuts against the paper's own method. The TOE constant is itself a
calorific-value convention, and if NEPRA reports on a gross basis rather than a net one, the
sector is overstated by \SIrange{5}{10}{\percent}. So the energy basis has not removed the
calorific-value problem; it has moved it from a per-fuel blend question to a single
documentation question about one constant. That is a better place for it---one page of one
yearbook resolves it---but it is not an escape, and Section~\ref{sec:limitations} treats it as
the sector's largest unresolved risk rather than a footnote.

\subsection{Coverage against accuracy}

Captive gas self-generation contributes \SI{1.25}{\percent} of national emissions and
\SI{0.03}{\percent} of national inventory variance. By the usual criterion for where inventory
effort should go---reduce the variance---it was not worth building.

It was worth building for a different reason, and the distinction is worth naming. This
inventory reports \SI{89}{\percent} of EDGAR's Pakistani sectors summed, and the shortfall is
not spread thinly across sectors but concentrated in categories that are entirely absent.
Absent categories contribute nothing to variance precisely because they contribute nothing at
all, so a variance-based priority ranking is blind to them by construction.

Coverage and accuracy are different objectives and need different criteria. The
key-category analysis tells you where to improve what you have; the completeness table tells
you what you do not have. Captive gas self-generation moved \SI{5.55}{\mega\tonne} from the second
column to the first, and no variance ranking would ever have recommended it.

\subsection{On labelling a bound as a bound}

The air-pollutant layer could have been published as an estimate. Applying a plausible removal
efficiency---say \SI{99}{\percent} for particulate and \SI{90}{\percent} for sulphur where
desulphurisation is installed---would have produced numbers that look like every other
published power-sector air inventory and would have been impossible for a reader to
distinguish from a measured result.

We publish the uncontrolled figure and label it, because the removal efficiency would have
been the largest single term in the calculation and entirely invented. The resulting number is
less immediately useful and more honest: it bounds what the fleet can emit, and it names the
control-device inventory as the specific missing input.

The same principle once produced a decision \emph{not} to report PM\textsubscript{2.5}: with
only a gas factor in hand, a gas-only figure would have been a real number, correctly
computed, and would have misled every reader who compared it against residential biomass. An
absent number is recoverable; a misleading present one is not.

That decision has since been reversed, and the reversal is worth recording rather than
quietly deleting. The blocker was never that AP-42 lacks a coal factor---it was that we had
not read far enough into the chapter. Table 1.1-6 gives the particle-size distribution for
bituminous and subbituminous combustion, Table 1.7-7 lignite's own, and applying each to
its rank's filterable mass yields
\SI{34.32}{\kilo\tonne} of PM\textsubscript{2.5} (Section~\ref{sec:air}). This is the same
not-absent-but-unread failure the transport and residential sectors each hit, and the same
correction: the honest response to a missing factor is to keep looking for it, and only then
to declare it missing.

\subsection{Implications}

\textbf{Read one page of the HDIP yearbook.} The TOE convention is worth
\SIrange{5}{10}{\percent} of this sector---more than the coal blend, more than the gas heating
value, more than any methodological refinement available---and resolving it requires a
conversion-factors table that exists and is paywalled rather than a measurement campaign.

\textbf{A control-device inventory would convert a bound into an estimate.} Which plants carry
precipitators and desulphurisation, and at what efficiency, is knowable from environmental
approvals and operator records. It is the difference between an air-pollutant layer that can
be cited and one that can only be bounded.

\textbf{Captive gas self-generation is a template for the rest of 1A2.} The method here---find the
regulator who regulates the \emph{fuel} rather than the plant, and reconcile against an
assembled register---closed \SI{13}{\percent} of the industrial-combustion gap with metered
data, and this inventory has since applied it to industrial process gas and to the industrial
liquid fuels metered by OCAC's Total Industry customer class, taking the covered share to
about a third. The remaining \SI{16}{\mega\tonne} is specifically solid: industrial
coal, petroleum coke and bagasse, in boilers, kilns and the captive plants that burn them
rather than gas. The same approach applied to those fuels' suppliers would close a
substantial further part without waiting for the paywalled yearbook.

\textbf{Provincial allocation deserves the same scrutiny as the total.} The capacity proxy
this paper replaced was wrong by \SI{7.9}{\mega\tonne} in one province and had been flagged as
probably wrong, in the right direction, in the parameter file. That it survived until
plant-wise generation became available says something about how rarely allocation proxies are
revisited relative to how often totals are.
